\section*{Results}\label{sec:results}

\subsection*{Category learning under binary and hierarchical feedback}
\label{r:task}

The model is formulated for three category-learning conditions using stimuli
defined by four normalized feature lengths. Condition~1 involves two
categories and binary feedback; condition~2 involves four categories and
binary feedback. Condition~3 also involves four categories, organized into
two families containing two species each. Feedback in this condition
distinguishes a correct species response, a response identifying the correct
family but the wrong species, and a response identifying the wrong family.
The learner is informed of the family--species structure and learns which
response labels belong to the same family. These conditions vary the
information available from a response and its feedback while retaining the
same model architecture.

\subsection*{Individual learning trajectories}
\label{r:behavior}
\draftslot{Results to be added.}

\subsection*{A dynamic rule model links hypothesis discovery to choice}
\label{r:model}

We formulated a dynamic rule model to describe how a learner discovers,
evaluates and uses candidate explanations of the category structure. The
model represents a rule as a partition of the stimulus space into labeled
category regions. A learner begins with one candidate and can subsequently
consider several rules, up to an individual capacity limit. Membership in
this candidate set determines which rules are available for prediction and
feedback learning. Each available rule has a relative belief that expresses
the support it receives within the current set, and a precision that
determines how concentrated its category predictions are.
These quantities jointly define a possible internal learning state.

Choices depend on both the perceived stimulus and the influence of the
available rules. Perception (P) introduces feature-specific uncertainty,
calibrated by independent perceptual measurements. Each candidate uses this
internal stimulus to predict categories from distances to its category
regions. Higher rule precision concentrates the prediction on nearby
regions. Choice readout (RC) combines these predictions using weights
derived from rule beliefs and an additional preference for the previous
trial's dominant rule. The dominant rule is the candidate with the largest
choice weight. Its continued influence allows an established response policy
to persist while evidence for alternatives accumulates. As beliefs and the
candidate set change, contributions to choice can become more dispersed or
more concentrated within the same participant (see \hyperref[m:choice]{Methods}).

Feedback changes both the relative support for rules and the precision of
their predictions. Evidence memory (M) combines the probability of the
received feedback under each candidate with the evidence retained from
preceding trials. Precision updating (B) uses the same predictions to
evaluate whether the feedback was more or less probable than under uniform
category predictions. Supporting feedback increases precision, whereas
contradictory feedback lowers it. All available candidates undergo these
updates, including candidates that contributed little to the response.
Belief and precision consequently describe different aspects of learning:
the former expresses support among competing explanations, and the latter
controls the concentration of a particular rule's predictions.

Hypothesis updating (H) changes the set of explanations considered on the
next trial. A state that summarizes recent failure controls the probability
of searching and the balance between local and global search. Local search
either edits a judgment in an existing rule or samples a rule with similar
classification outcomes. Global search samples from the remaining rule
catalogue. Rules with lower relative belief have a higher probability of
leaving the candidate set. The number of candidates can therefore grow as
new hypotheses enter and shrink as weak hypotheses are abandoned, within
the participant's capacity limit. A new candidate receives a share of each
retained rule's belief, with larger shares transferred from rules with more
similar classification outcomes. This operation expresses the
generalization of existing classification knowledge to a newly considered
explanation (see \hyperref[m:hypothesis]{Methods}).

The feedback conditions also determine which latent relations are learned.
In the hierarchical condition, the model maintains joint beliefs over
classification rules and the three possible pairings of response labels
into families. Fully correct, partially correct and incorrect responses
provide evidence about this joint state. When a new rule enters, it inherits
the accumulated marginal belief about response pairings. This preserves
learned family relations as the candidate classification rules change
(see \hyperref[m:pairing]{Methods}).

The model thus defines several points at which learning can change. A useful
rule can enter the candidate set before it receives substantial belief;
increasing belief can precede its becoming dominant; and its predictions can
become more precise while it remains dominant. Candidate-set size measures
how many explanations are available, whereas effective participation
measures how concentrated their contributions to choice are. These states
provide separate quantities for examining rule discovery, belief
accumulation, changes in rule use and increasing prediction precision
(see \hyperref[m:measures]{Methods}).

Participant-level parameters govern capacity, evidence retention, search,
choice inertia and precision learning. They remain fixed within a
participant, while the candidate set and its associated states evolve across
trials. Particle filtering approximates the distribution of internal paths
consistent with the available choice and feedback history. Predictions are
averaged using path weights determined before the current choice and its
feedback, defining the choice likelihood. Weighted states describe
uncertainty about the inferred learning process
(see \hyperref[m:inference]{Methods}).

\subsection*{Behavioral prediction and model comparison}
\label{r:prediction}
\draftslot{Results to be added.}

\subsection*{Latent trajectories of hypothesis discovery and use}
\label{r:trajectories}
\draftslot{Results to be added.}

\subsection*{Effects of feedback structure on learning}
\label{r:feedback}
\draftslot{Results to be added.}

\subsection*{Generative predictions and model identifiability}
\label{r:recovery}
\draftslot{Results to be added.}
